\chapter{Oxygen Therapy}

\section*{Overview}
Oxygen therapy provides supplemental oxygen to the lungs to treat hypoxemia (low blood oxygen levels). It increases alveolar oxygen tension and improves tissue oxygenation.

\section*{Alternative Names}
Supplemental oxygen, O2 therapy

\section*{Principle}
Supplemental oxygen is delivered by mask, nasal cannula, or other devices at concentrations above ambient air, increasing the fraction of inspired oxygen. This raises arterial oxygen saturation and improves oxygen delivery to tissues when normal breathing is insufficient.
\section*{History}
Joseph Priestley discovered oxygen in 1774. Alvan Barach developed modern oxygen tents, masks, and nasal cannulas in the early 20th century.

\section*{Why It Is Done}
Ordered for acute respiratory failure, COPD exacerbation, pneumonia, heart failure, sleep apnea, or carbon monoxide poisoning.

\section*{Is This a Primary or Secondary Test or Procedure}
Primary supportive therapy for patients with acute or chronic hypoxemic respiratory failure.

\section*{How It Is Performed}
Oxygen is delivered via a nasal cannula (prongs in the nose), a face mask, or high-flow nasal therapy. The flow rate is adjusted based on pulse oximetry.

\section*{Preparation}
No preparation. Arterial blood gas (ABG) may be performed to establish baseline oxygenation.

\section*{Contraindications and Precautions}
Oxygen is prescribed when the expected benefit outweighs the risks, with a target saturation set for the individual. Many acutely ill adults are treated toward 94--98\%; people at risk of hypercapnic respiratory failure may require a lower target, often 88--92\%, with blood-gas or other monitoring as directed. Do not adjust flow or discontinue oxygen without clinical advice.

\section*{Risks}
Oxygen toxicity, carbon dioxide retention, fire hazard, or nasal dryness and irritation.

\section*{Aftercare and Recovery}
Monitor oxygen saturation and work of breathing. Adjust flow rate to maintain target saturation (typically 92-96%, or 88-92% in COPD).

\section*{Outcome and Follow-up}
Monitored using pulse oximetry (SpO2) and arterial blood gases. Resolution of hypoxemia is the primary goal.

\section*{Expected Results}
Maintenance of target arterial oxygen saturation (SpO2 92-96% or 88-92% depending on the clinical population).

\section*{Follow-up}
Pulse oximetry checks, titration of oxygen flow, or pulmonary consult.

\section*{When to Seek Medical Care}
Follow the procedure team's specific instructions. Seek urgent medical assessment for severe or worsening pain, uncontrolled bleeding, fever or signs of infection, trouble breathing, chest pain, fainting, new weakness or confusion, inability to urinate, or any rapidly worsening symptom. The appropriate warning signs depend on the procedure and the patient's medical condition.

\section*{References}
\begin{enumerate}
  \item MedlinePlus. Oxygen Therapy. U.S. National Library of Medicine; 2025.
  \item Current Medical Diagnosis and Treatment. McGraw-Hill; 2025.
  \item British Thoracic Society. Emergency oxygen guideline. \url{https://www.brit-thoracic.org.uk/clinical-resources/guidelines/emergency-oxygen/} (accessed 2026-08-16).
\end{enumerate}

\clearpage
